\documentclass[10pt,letterpaper]{article}
\usepackage{cornell-notes}

\title{Clause 29.07.10.04: Conversions Between System Clock And Other Clocks}
\author{}
\date{}
\setDocTitle{Clause 29.07.10.04: Conversions Between System Clock And Other Clocks}
\setDocSubtitle{C++ 2024}
\setDocOwner{C++ 2024 Cornell Notes}
\setDocDate{\today}
\setCornellCollection{C++ 2024 Cornell Notes}
\setCornellUnitType{Clause}
\setCornellUnitNumber{29.07.10.04}
\setCornellUnitTitle{Conversions Between System Clock And Other Clocks}
\hypersetup{pdftitle={Clause 29.07.10.04: Conversions Between System Clock And Other Clocks},pdfsubject={Cornell notes},pdfauthor={}}

\begin{document}
\maketitle
\begin{CornellOverviewBox}{Overview}
\textbf{Classification:} Standard library - Normative\\[2pt]
\textbf{Learning objective:} Understand the rules, terminology, and semantic consequences associated with Conversions between system\_clock and other clocks.
\end{CornellOverviewBox}\begin{CornellSummaryBox}{Summary}
{\sffamily\bfseries\color{Accent} Summary}\par\vspace{3pt}Section 29.7.10.4 focuses on Conversions between system\_clock and other clocks. Apply the specified interface together with its mandates, effects, result conditions, exception behavior, and complexity constraints. When applying it, preserve the distinction between mandatory requirements, recommendations, permissions, and behavior deliberately left undefined, unspecified, or implementation-defined.
\end{CornellSummaryBox}

\end{document}
